%%%%%%%%%%%%%%%%%%%%%%%%%%%%%%%%%%%%%%%%%%%%%%%%%%%%%%%%%%%%%%%%%%%%%%%%%%%%%%%
%% Capítulo 3 --- Desenvolvimento e métodos
%%
%% Descreve os experimentos e o que está em kaelix-firmware: componentes reais,
%% firmware, receptor e ensaio de bancada. O projeto de hardware (placas,
%% invólucro, térmica, bateria) é trabalho futuro, por decisão do usuário.
%%
%% Valores conferidos em platformio.ini, src/sensors/vibration.h,
%% src/comms/lora_config.h, src/comms/lora.h, src/main.cpp e logs/.
%%%%%%%%%%%%%%%%%%%%%%%%%%%%%%%%%%%%%%%%%%%%%%%%%%%%%%%%%%%%%%%%%%%%%%%%%%%%%%%

\makeatletter
\@ifundefined{toprule}{%
  \newcommand{\toprule}{\hline}%
  \newcommand{\midrule}{\hline}%
  \newcommand{\bottomrule}{\hline}%
}{}
\@ifundefined{emandamento}{%
  \@ifundefined{textcolor}%
    {\newcommand{\emandamento}[1]{#1}}%
    {\newcommand{\emandamento}[1]{\textcolor{red}{#1}}}%
}{}
\makeatother

\chapter{Desenvolvimento e métodos}
\label{cap:desenvolvimento}

O trabalho tem duas partes. A primeira são os experimentos computacionais, que usam sinal
sintetizado e gravações públicas de bancada para avaliar a aquisição e a estimação. A
segunda é o dispositivo: um acelerômetro e um rádio LoRa ligados a um microcontrolador
ESP32, o \emph{firmware} que os controla e o receptor.

\section{Delimitação}
\label{sec:dev-delimitacao}

O objeto de estudo é o Kaelix, um dispositivo de monitoramento de vibração de baixo custo.
Este trabalho trata da cadeia de medição, do \emph{firmware} e da comunicação. O projeto de
placas de circuito impresso, o invólucro, a bateria e o comportamento térmico ficam para
trabalho futuro (seção~\ref{sec:conc-futuro}). Nos ensaios, os componentes foram ligados
em montagem de bancada e alimentados pela porta USB.

A norma ISO 20816-3 serve aqui como referência para comparar métodos de integração. Ela
exige resposta plana de \SIrange{10}{1000}{\hertz}, e o Kaelix não atende a esse requisito,
porque o filtro interno do sensor limita a banda a \SI{184}{\hertz}. Por isso, o
dispositivo não classifica máquinas nas zonas da norma.

\section{Modelo de sinal e cadeias de aquisição}
\label{sec:dev-sinal}

O efeito da aquisição foi medido sobre dois tipos de sinal: um sinal sintetizado, cuja
resposta correta é conhecida, e gravações reais decimadas para a taxa do dispositivo.

\subsection{Sinal sintetizado}

O sinal sintetizado combina dois mecanismos. Desbalanceamento e desalinhamento aparecem como
harmônicas da rotação, todas abaixo de \SI{100}{\hertz}. O defeito de pista externa aparece
como um trem de impulsos amortecidos, repetido na taxa característica da geometria, que
modula uma ressonância estrutural. A aceleração simulada é

\begin{equation}
a(t) \;=\; \underbrace{\sum_{k=1}^{3} A_k \sin\!\left(2\pi k f_r t + \varphi_k\right)}_{\text{banda baixa}}
\;+\; \underbrace{\sum_{i} D\, e^{-\zeta_d (t-t_i)} \sin\!\left(2\pi f_{\mathrm{res}} (t-t_i)\right) u(t-t_i)}_{\text{impactos do defeito}}
\;+\; \varepsilon(t),
\label{eq:dev-sinal}
\end{equation}

\noindent
com $f_r$ a rotação, $A_k$ as amplitudes harmônicas, $D$ a amplitude do defeito, $\zeta_d$
a taxa de amortecimento do impacto, $f_{\mathrm{res}}$ a ressonância excitada,
$u(\cdot)$ o degrau unitário e $\varepsilon(t)$ ruído gaussiano branco. Os instantes de
impacto são $t_i = i/\mathrm{BPFO} + \delta_i$, em que $\delta_i$ modela o escorregamento
das esferas. As frequências características vêm da geometria do rolamento, pelas expressões
da seção~\ref{sec:fund-rolamento}.

\subsection{Três cadeias comparadas sobre o mesmo sinal}

O mesmo sinal passou por três cadeias:

\begin{alineas}
  \item referência: amostragem a \SI{20}{\kilo\hertz}, que representa um acelerômetro de
        banda larga;
  \item com filtro: decimação para \SI{1}{\kilo\hertz} com filtro antirrebatimento;
  \item sem filtro: subamostragem para \SI{1}{\kilo\hertz} sem nenhum filtro.
\end{alineas}

No terceiro caso, uma componente acima da frequência de Nyquist é observada na frequência
aparente da Equação~\eqref{eq:fund-alias}.

Nenhuma das cadeias simuladas é idêntica à do dispositivo, que usa o filtro interno do
sensor com banda de \SI{184}{\hertz}. A segunda cadeia mantém conteúdo até
\SI{500}{\hertz}. A terceira é um caso extremo, usado para medir o rebatimento isolado,
porque o MPU6050 sempre opera com o filtro interno (Tabela~\ref{tab:fund-dlpf}).
\emandamento{Os experimentos ainda precisam ser refeitos com a banda de \SI{184}{\hertz}.}

Na comparação entre cadeias, os descritores foram a curtose, o fator de crista, o pico e o pico
a pico, calculados sobre \SI{2}{\second} de sinal. Os experimentos de limiar, métrica e partição
usam janelas de \SI{0,2}{\second}, com \SI{50}{\percent} de sobreposição, e outro conjunto: valor
eficaz, curtose, fator de crista, média absoluta, razão entre valor eficaz e média absoluta e fração
da energia por banda. A cadeia de \SI{20}{\kilo\hertz} tem quatro bandas e soma nove descritores, e
a de \SI{1}{\kilo\hertz} tem duas e soma sete.

O detector do \emph{firmware} usa quatro descritores: valor eficaz, curtose em excesso, fator de
crista e frequência dominante do espectro de velocidade. As fórmulas são iguais no treino e no
\emph{firmware}, com duas diferenças de sinal: o \emph{firmware} remove a média da janela e lê o eixo
Z, e o treino não remove a média e usa o canal radial do sensor \emph{underhang}. Nenhum desses
descritores mede a relação entre harmônicas, a fase ou a linha na frequência característica
(seção~\ref{sec:res-discussao}).

\subsection{Dado real decimado}

As taxas de detecção por modo de falha e a comparação entre limiares vieram das gravações do conjunto público
MAFAULDA~\cite{mafaulda}, obtidas em bancada instrumentada a \SI{50}{\kilo\hertz} e
decimadas em \emph{software} para \SI{1}{\kilo\hertz}. O conjunto soma \num{1951} arquivos
em seis classes (operação normal, desbalanceamento, desalinhamento horizontal e vertical e
defeito de rolamento nas posições \emph{overhang} e \emph{underhang}), segmentados em
\num{17559} janelas de \num{512} amostras, sem sobreposição. Apenas \num{49} arquivos são de operação normal,
e todo o treino saiu deles, porque o detector aprende somente a operação normal. A bancada
está abaixo do limite de potência e de altura de eixo da norma. Ela serve como fonte de
sinal para estudar a cadeia de medição, mas não como caso de aplicação da norma.

O limiar do detector foi calibrado pelo quantil do escore sobre arquivos normais não usados no
ajuste das árvores, com falso alarme alvo de \SI{1}{\percent}. O limiar padrão de \num{0,5} foi
avaliado nas mesmas partições.

A validação usou cinco partições por ensaio, com os arquivos de origem como grupos, para
que janelas de um mesmo ensaio não ficassem no treino e no teste ao mesmo tempo
(seção~\ref{sec:res-estimacao}). Nos ensaios de rolamento, os resultados foram separados
pelo nível de massa de desbalanceamento sobreposta, porque o valor agregado mistura os dois
efeitos.

\section{Conversão normativa e detecção}
\label{sec:dev-estimacao}

A norma define as zonas de severidade em velocidade eficaz de banda larga, e o acelerômetro
entrega aceleração \cite{iso208163_2022}. A conversão foi testada contra uma senoide cuja
integral é conhecida analiticamente, em cinco variantes: integração no tempo com e sem
filtro passa-alta, integração no domínio da frequência restrita à banda normativa, e esta
última e a integração no tempo sem filtro repetidas sob viés contínuo contido na
especificação do sensor. O método correto já está na literatura \cite{brandt2014}, e este
trabalho mediu o erro das alternativas em relação ao valor analítico. O ensaio usa
\SI{5}{\kilo\hertz} e \SI{4}{\second} de sinal. A integração no tempo é uma soma acumulada que parte
de zero, sem remoção da média, e o filtro passa-alta, quando usado, é de \SI{10}{\hertz}.

A detecção usou o Isolation Forest \cite{liu2008}, treinado apenas com dados de operação normal,
como no caso de uma máquina sem histórico rotulado. Em sinal sintetizado, na cadeia de
\SI{20}{\kilo\hertz}, o parâmetro de contaminação variou em cinco valores. O falso alarme foi medido
nas mesmas janelas do ajuste e, depois, em um ensaio normal deixado fora do treino, em cinco rodadas.

A métrica foi avaliada em defeito incipiente, com amplitude $D$ de \num{0,25}, contra \num{2,0} no
experimento de limiar. A área sob a curva ROC global e a área parcial \cite{koizumi2020} foram
comparadas entre as cadeias de \SI{20}{\kilo\hertz} e de \SI{1}{\kilo\hertz}, sobre as janelas
normais do ajuste e as janelas com defeito.

O protocolo de validação foi avaliado com um Random Forest supervisionado de \num{150} árvores,
treinado com as duas classes, na cadeia de \SI{1}{\kilo\hertz}. Esse classificador não é o detector
do \emph{firmware}.

\section{Dispositivo e \emph{firmware}}
\label{sec:dev-firmware}

O dispositivo é formado por um acelerômetro MEMS ligado por I2C a um microcontrolador ESP32
e por um módulo de rádio RA-02, com o circuito SX1278, ligado por SPI. A
Tabela~\ref{tab:dev-config} mostra a configuração usada nos ensaios. \emandamento{O sensor
da montagem de 24 de outubro de 2026 respondeu ao registro de identidade com o valor
\texttt{0x74}, do perfil MPU6515, e não com o \texttt{0x68} do MPU6050. A marcação do chip
ainda será conferida. As Tabelas~\ref{tab:fund-mpu6050} e~\ref{tab:fund-dlpf} são do MPU6050.}

\begin{table}[!htb]
\IBGEtab{%
\caption{Configuração do dispositivo nos ensaios de bancada}%
\label{tab:dev-config}%
}{%
\begin{tabular}{@{}ll@{}}
\toprule
Item & Valor \\
\midrule
Microcontrolador         & ESP32-S3 ou ESP32-C3 \\
Faixa do acelerômetro    & $\pm\SI{4}{g}$ \\
Filtro interno do sensor & \texttt{DLPF\_CFG}~=~1, banda de \SI{184}{\hertz} \\
Taxa de amostragem alvo  & \SI{1}{\kilo\hertz}, com tolerância de \SI{2}{\percent} \\
Janela                   & \num{512} amostras do eixo Z \\
Frequência do rádio      & \SI{433}{\mega\hertz} \\
Parâmetros do enlace     & SF9, BW de \SI{125}{\kilo\hertz}, CR de 4/5, preâmbulo de 8 símbolos \\
Potência de transmissão  & \SI{2}{dBm} \\
Pacote                   & \num{21} bytes, com CRC de 16 bits \\
Intervalo entre envios   & cerca de \SI{11}{\second} \\
\bottomrule
\end{tabular}%
}{%
\fonte{Elaboração própria.}%
}
\end{table}

O \emph{firmware} foi escrito em C++ com o \emph{framework} Arduino e compilado com o
PlatformIO. O mesmo código atende às duas placas, com a pinagem definida em um arquivo
próprio. Na bancada, cada
ciclo configura o acelerômetro, lê a janela pelo FIFO do sensor, calcula os quatro descritores,
envia o pacote pelo rádio e espera cerca de \SI{10}{\second} até o ciclo seguinte. O pacote leva
o identificador do dispositivo, o contador de ciclos, o estado, o código de erro, o valor eficaz e a
temperatura, que vai como NaN na bancada. Os outros três descritores só aparecem no registro serial.

Durante a leitura, o \emph{firmware} estima a taxa de amostragem real e descarta a janela
se ela se afastar mais de \SI{2}{\percent} de \SI{1}{\kilo\hertz}, porque a frequência
dominante é calculada com base nessa taxa.

Cada função devolve um código de estado, e o pacote leva o código do primeiro erro do ciclo.
Quando a leitura falha, a função não escreve no vetor de saída, o pacote sai com o estado
\texttt{UNKNOWN} e o valor eficaz vai como NaN. Uma versão anterior preenchia o vetor com
zeros em caso de erro, e os descritores calculados sobre esses zeros seriam classificados
como máquina sadia. O Apêndice~\ref{ape:praticas} descreve essas práticas.

O detector tem \num{100} árvores, treinadas em Python com subamostras de \num{256} pontos e
exportadas para um arquivo de cabeçalho em C++. Um teste compara o escore calculado em C++
com o da biblioteca usada no treino. \emandamento{A inferência ainda não é executada na
bancada: nesse perfil, o \emph{firmware} envia sempre o estado \texttt{UNKNOWN}.}

As funções de processamento não dependem do \emph{hardware} e rodam no computador. Os
\num{85} casos de teste da suíte passaram na versão avaliada, e parte deles cobre módulos fora da
montagem, como o termistor.

\section{Receptor}
\label{sec:dev-comunicacao}

O receptor é outro ESP32 com o mesmo módulo de rádio. Os parâmetros do enlace ficam em um
único arquivo, compartilhado com o transmissor, para que uma atualização da biblioteca de
rádio não os altere. Antes de aceitar um pacote, o receptor confere o tamanho, o CRC, a
versão e a validade dos campos. O dispositivo não tem relógio e envia apenas um contador de
ciclos, então é o receptor que informa o instante de chegada ao enviar os dados. A diferença entre
contadores consecutivos mostra quantos pacotes se perderam.

A implementação do CRC em C++ foi comparada com uma implementação em Python em \num{204} casos
de teste. \emandamento{O envio dos dados do receptor a um servidor por
Wi-Fi está implementado, mas não tem registro de ensaio.}

\section{Ensaio de bancada}
\label{sec:dev-bancada}

\emandamento{Foram analisados dois registros de bancada, salvos no repositório do projeto.
No primeiro, de 11 de outubro de 2026, um ESP32-S3 transmitiu e um ESP32-C3 recebeu, e o
registro é o do receptor. No segundo, de 24 de outubro de 2026, um ESP32-C3 transmitiu, e o
registro é o do transmissor, sem receptor associado. Os dois foram usados para avaliar os
critérios da hipótese H6: o desvio da taxa de amostragem em relação ao limite de
\SI{2}{\percent} e a validação dos pacotes pelo CRC. Esses registros serão substituídos
pelas medições com o dispositivo montado.}
